\documentclass[10pt,a4paper]{amsart}
\usepackage[margin=19mm]{geometry}
\usepackage{amsmath,amssymb,mathtools}
\usepackage[scr=rsfs,cal=euler]{mathalfa}
\usepackage{xfrac}
\usepackage[unicode,hidelinks,bookmarks=false]{hyperref}
\newcommand{\ie}{i.e.\ }
\newcommand{\cf}{cf.\ }
\newcommand{\resp}{respectively\ }
\newcommand{\Subsection}{\subsection}
\providecommand{\nobreakdash}{\nobreak\hbox{-}}
\setlength{\emergencystretch}{2em}
\multlinegap0pt
\usepackage{braket}
\usepackage{amssymb}
\usepackage{mathtools}
\usepackage{bm}
\usepackage[shortlabels]{enumitem}
\newtheorem{proposition}{Proposition}
  \def\Jac{J}%
  \def\bra#1{<#1|}%
  \def\ee{e}%
  \def\ii{i}%
  \def\ket#1{|#1>}%
\title{Polynomial Initial-State Jumps and Christoffel Transforms in Krylov Complexity}
\author{Abhishek Chowdhury, Ajit Prasad Mahapatra}
\date{Source: arXiv:2607.05294v3 (CC BY 4.0)}
\begin{document}
\maketitle

\noindent\emph{Source and licence.} Polynomial Initial-State Jumps and Christoffel Transforms in Krylov Complexity. Version arXiv:2607.05294v3 (CC BY 4.0), distributed by arXiv under the Creative Commons Attribution 4.0 International licence, http://creativecommons.org/licenses/by/4.0/, as recorded in the arXiv licence metadata for this exact version and shown on the abstract page https://arxiv.org/abs/2607.05294v3. Full licence notice: \texttt{licenses/arxiv-2607.05294-CC-BY-4.0.txt}.\par\smallskip

\noindent\textit{Editorial scope.} Counted unit: the article's proposition prop:chord-radial-filtration (source0.tex lines 3654--3678). The article's complete original proof of it is delivered with it (source0.tex lines 3680--3711, one \texttt{proof} environment of 32 lines). No mathematics is written, repaired or re-derived: the statement and the proof are the article's own lines, in the article's own notation, and no retained line is substituted or re-flowed.  No further source-local context is needed: every cross-reference of every retained line resolves inside the two delivered spans.  Nothing was omitted from any retained span, so the package carries no omission record.  No retained line contains a citation, so the wrapper carries no bibliography and no bibliography entry.  No retained line contains a figure environment, an image-inclusion command or drawing code, so no image is delivered and the package declares no asset dependency.  Editorial formatting only, in addition: an AMS \texttt{amsart} preamble that restates the article's own declarations and macros used by the retained lines, namely the environments \texttt{proposition} on separate counters, together with the standard packages those lines need.  The retained environments print in this document as Proposition 1 in order of appearance, whereas the article prints the same items with the numbering of its own full document.  The complete native source of the article as distributed in the arXiv e-print for this exact version is bundled unchanged as \texttt{sources/204530/source0.tex}.
\begin{proposition}[Radial filtration for a fixed chord]
\label{prop:chord-radial-filtration}
For every $r\ge0$, $0\le q<1$ and finite $t$, let
\begin{equation}
   B_{r,\ell}=\sum_{\substack{m\ge0\\ |m-r|\le\ell}}
      \ket m_{\Jac}{}_{\Jac}\!\bra m,
   \qquad
   \widetilde\Pi_{\ell}^{(r)}=
      \sum_{n=0}^{\ell}\ket{\widetilde K_n^{(r)}}
      \bra{\widetilde K_n^{(r)}}.
   \label{eq:chord-radial-projectors}
\end{equation}
Then
\begin{equation}
\begin{aligned}
   K_r(t)-D_r^{\rm ch}(t)
   &=\sum_{\ell\ge0}\bra{\psi_r(t)}
     \bigl(B_{r,\ell}-\widetilde\Pi_{\ell}^{(r)}\bigr)
     \ket{\psi_r(t)}\ge0\,,\\
   K_r(t)&\ge D_r^{\rm ch}(t)\ge
      \bigl|\overline n_r(t)-r\bigr|.
   \label{eq:chord-radial-filtration}
\end{aligned}
\end{equation}
\end{proposition}

\begin{proof}
Tridiagonality implies that the first $\ell+1$ vectors generated from
$\ket r_{\Jac}$ have physical chord support only at distances at most $\ell$.
Consequently,
$\operatorname{Ran}\widetilde\Pi_{\ell}^{(r)}\subseteq
\operatorname{Ran}B_{r,\ell}$.  The two orthogonal projectors are nested, so
$B_{r,\ell}-\widetilde\Pi_{\ell}^{(r)}$ is itself an orthogonal projector.
For a nonnegative integer-valued random variable, its first moment is the sum of its
tail probabilities.  Applied in the two bases, this identity gives
\begin{equation}
\begin{aligned}
 K_r(t)&=\sum_{\ell\ge0}
 \left[1-\bra{\psi_r(t)}\widetilde\Pi_{\ell}^{(r)}\ket{\psi_r(t)}\right],\\
 D_r^{\rm ch}(t)&=\sum_{\ell\ge0}
 \left[1-\bra{\psi_r(t)}B_{r,\ell}\ket{\psi_r(t)}\right],
\end{aligned}
\end{equation}
where these sums converge.  Indeed,
$\|\Jac_q\|\le R_q$, where $R_q=2/\sqrt{1-q}$, and a matrix element between chord
levels separated by $d$ vanishes in every power $\Jac_q^k$ with $k<d$.
Hence, for $|t|\le T$,
\begin{equation}
 \left|{}_{\Jac}\!\bra m\ee^{-\ii t\Jac_q}\ket r_{\Jac}\right|
 \le \sum_{k=|m-r|}^{\infty}\frac{(R_qT)^k}{k!}\,.
\end{equation}
The same estimate holds in the rebuilt Jacobi representation, whose spectral support is
contained in that of $\mu_q$.  The squares of these bounds are summable with the respective
distance or level weights.  Subtracting the two tail sums proves the first line of
eq.~\eqref{eq:chord-radial-filtration}.  Finally,
$D_r^{\rm ch}(t)\ge|\overline n_r(t)-r|$ is the triangle inequality for the
physical chord-number distribution.
\end{proof}

\end{document}
